%% kapa: exam=1 round=28 year=2017 period=2 form=A subject=appraisal_law subject_ko=감정평가관계법규 seq=8 number=8
%% source: Q-Net 감정평가사 기출문제 (cst00309) artlSeq=5199603 file=extracted/1008/5199603/2017년 제28회 감정평가사 1차시험 문제지/2교시 A형.pdf
%% conversion: kordoc 4.24.51 -> tools/gen1.py (auto); review=proofread-sonnet; flags=none
\begin{question}{8}
국토의 계획 및 이용에 관한 법령상 공동구에 관한 설명으로 옳지 않은 것은?
\begin{choices}
  \item 200만 제곱미터를 초과하는 택지개발지구에서 개발사업을 시행하는 자는 공동구를 설치하여야 한다.
  \item 공동구가 설치된 경우 가스관은 공동구협의회의 심의를 거쳐 공동구에 수용할 수 있다.
  \item 공동구의 설치에 필요한 비용은 공동구 점용예정자가 부담하되, 그 부담액은 사업 시행자와 협의하여 정한다.
  \item 공동구의 효율적인 관리ㆍ운영을 위하여 필요하다고 인정하는 경우에는 지방공사 또는 지방공단에 그 관리ㆍ운영을 위탁할 수 있다.
  \item 공동구관리자는 공동구의 안전점검으로서 매년 1월 1일을 기준으로 6개월에 1회 이상 정기점검을 실시하여야 한다.
\end{choices}
\end{question}
\answer{3}
